\section{Conclusions}
We finally arrive at conclusions and what we've concluded is that there
is no Unified Field Theory, there is not even a Grand Unified Theory, and the place to look for
a deeper understanding of the realtionship between quantum mechanics and general relevitiy is
not at the Planck scale but at the scale where records are confirmed in a network of events.

Although there is no Unified Field Theory, there is a unifying principle, an ungrounded epistemological self-referential state of knowledge made
to appear as if it is a non-self-referential ontological causal sequence as suggested in
\autocite{small_2006}.
Using this principle we have been able to make a start at developing a rationale for the
Standard Model and for the 3+1 structure of spacetime. 

Now, I cannot claim that idea as my own, I learned it from Indian philosopher Maharish Mahesh Yogi. My 
only contribution has been to use that principle as a thread to stich together many ideas from other 
people into a coherent tapestry.

One thing I am very uncertain about is the issue of mass. We've settled on mass as \enquote{associator debt}
and then discovered the idea has a better foundation as \enquote{norm defect} in the sedenions. Which sort
of makes sense since the sedenions give rise to the 3 generations of fermions, which are distinguished by
their masses. I think that's a temporary proxy for a deeper idea. We haven't linked mass to the causal network we 
experience as spacetime and I don't think we have the right idea until that link can be made.

Since we're talking about self-reference, then hanging in the background is the Lie group E8,
the only Lie group that acts on itself. Along with developing the idea of a computational
network and attempting to identify it with real space and time I'm postponing that investigation
to paper2. A very rough outline is present in the Claude brain dump
\autocite{small2026itfrombit}.
\subsection{Updates}
While doing the final preparations for release I realised there's a mistable in the idea that
deletion post-FANOUT has a thermodynamic cost and pre-FANOUT it doesn't. I think deletion 
pre-FANOUT has a complex thermodynamic cost. I'll be exploring this idea, but it's too
late to make the cut for this version
